% Table: A-level subject combinations by economics availability (Chapter 4 appendix)
\begin{table}[htbp]
\centering
\caption{A-level subject combinations by economics availability}
\label{tab:subject_triples}
\begin{tabular}{l cc cc}
\toprule
                                  & \multicolumn{2}{c}{Switcher schools}  & Always     & Never      \\
\cmidrule(l{3pt}r{3pt}){2-3}
                                  & Econ        & Econ not    & offer      & offer      \\
                                  & available   & available   &            &            \\
\midrule
\multicolumn{5}{c}{\textbf{Panel A: Most common subject triples (\% of students)}} \\
\midrule
Maths, Economics, Physics         & 0.32          & ---         & 0.46         & ---        \\
Maths, Business, Physics          & 0.20          & 0.25          & 0.17         & 0.23         \\
Maths, Biology, Chemistry          & 3.29          & 3.08          & 3.86         & 2.91         \\
Maths, Geography, Physics         & 0.34          & 0.46          & 0.34         & 0.40         \\
Maths, History, Physics           & 0.21          & 0.21          & 0.17         & 0.21         \\
Sciences triple                   & 0.94          & 0.88          & 0.73         & 1.01         \\
Humanities triple                 & 10.2          & 11.2          & 8.4         & 11.1         \\
Other combinations                & 84.5          & 84.0          & 85.9         & 84.1       \\
\midrule
\multicolumn{5}{c}{\textbf{Panel B: Compliers' counterfactual subject bundle (\%)}} \\
\midrule
Business/management               & \multicolumn{2}{c}{24.7}                &            &            \\
Mathematics                       & \multicolumn{2}{c}{-7.7}                &            &            \\
Geography                         & \multicolumn{2}{c}{1.5}                &            &            \\
History                           & \multicolumn{2}{c}{8.0}                &            &            \\
Sciences                          & \multicolumn{2}{c}{-1.8}                &            &            \\
English literature                          & \multicolumn{2}{c}{23.3}                &            &            \\
Psychology                          & \multicolumn{2}{c}{7.0}                &            &            \\
\addlinespace
Observations                      & 104,374          & 117,200          & 321,263         & 138,754         \\
\bottomrule
\end{tabular}
\begin{minipage}{\textwidth}
\smallskip
\footnotesize
\textit{Notes:} Panel~A compares the distribution of A-level subject triples across school-cohorts in switcher schools when economics is and is not available, and in always-offer and never-offer schools. Panel~B reports the composition of the subjects compliers would have taken instead of economics. Each subject's take-up indicator is regressed on economics availability within switcher schools, with school and cohort fixed effects; every student in this sample takes exactly three A-levels, so each student who takes up economics drops exactly one other subject. The entry is the fall in that subject's take-up divided by the rise in economics take-up, expressed as a percentage. It is the share of compliers who would otherwise have taken that subject. Shares sum to 100 across all subjects; those shown account for 55\%, with the remainder spread across smaller subjects. A negative entry means the subject is taken more, not less, when economics is available: mathematics and the sciences are complements to economics rather than substitutes. The counterfactual mix in Section~\ref{sec:lifetime} is weighted by the positive entries. Sample: GCSE cohorts 2002--2008, students taking exactly three A-levels.
\end{minipage}
\end{table}
